{\color{unimain}\section[Пример расчета уставок защит трансформатора]{Пример расчета уставок защит трансформатора}}
\color{black}

\color{uniblue}\subsection[Исходные данные]{Исходные данные} 
\color{black}

\begin{enumerate}
\item
В настоящем примере дан расчет уставок защит и автоматики понижающего двухобмоточного трансформатора типа ТМН-4000/35 с параметрами, приведенными в таблице \ref{sett:tbl1}. Поясняющая схема приведена на рисунке \ref{sett:img1}.

\vspace{3mm}
\begin{figure}[H]
\centering
\includegraphics[width=0.5\textwidth,height=0.5\textheight,keepaspectratio]{\genpath/000. Общее/schemasT.pdf}
\captionof{figure}{Поясняющая схема и схема замещения защищаемого трансформатора}\label{sett:img1}
\end{figure}

\small
\begin{longtable}{|>{\centering\arraybackslash}m{7.2cm}|>{\centering\arraybackslash}m{3cm}|>{\centering\arraybackslash}m{3cm}|>{\centering\arraybackslash}m{3cm}|}
\caption{Параметры защищаемого трансформатора\hfill\vspace{-0.5\baselineskip}}\label{sett:tbl1}\\ 
\hline
\rowcolor{gray!30}
Параметр & Обозначение параметра & Единица измерения & Значение \\ 
\hline
\endfirsthead
\caption*{\hspace{3pt}\emph{Продолжение таблицы \ref{sett:tbl1}\hfill\vspace{-0.5\baselineskip}}} \\ % сделано по ГОСТ 2.105 п.6.8.7
\hline
\rowcolor{gray!30}
Параметр & Обозначение параметра & Единица измерения & Значение \\ 
\endhead
\endfoot
\endlastfoot
\raggedright 1 Схема соединения & \centering -- & \centering -- & \centering\arraybackslash Y/$\Delta$-11 \\
\hline
\raggedright 2 Номинальная мощность & \centering Sном & \centering МВА & \centering\arraybackslash 4 \\
\hline
\raggedright 3 Номинальное напряжение обмотки ВН & \centering Uном,в & \centering кВ & \centering\arraybackslash 35 \\
\hline
\raggedright 4 Номинальное напряжение обмотки НН & \centering Uном,н & \centering кВ & \centering\arraybackslash 6,3 \\
\hline
\raggedright 5 Диапазон регулирования напряжения РПН со стороны ВН & \centering -- & \centering -- & \centering\arraybackslash \pm 6x1,5\% \\
\hline
\raggedright 6 Напряжение короткого замыкания & \centering Uк & \centering \% & \centering\arraybackslash 7,5 \\
\hline
\end{longtable}
\normalsize

\item 
Схема замещения и расчетные значения токов короткого замыкания приведены на рисунке \ref{sett:img2}.

\vspace{3mm}
\begin{figure}[H]
\centering
\includegraphics[width=1\textwidth,height=1\textheight,keepaspectratio]{\genpath/000. Общее/schemasT2.pdf}
\captionof{figure}{Схема замещения защищаемого трансформатора и результаты расчета токов КЗ}\label{sett:img2}
\end{figure}

\item
На трансформаторе установлены трансформаторы тока на стороне ВН с коэффициентом трансформации 400/5, на стороне НН с коэффициентом трансформации 1500/5. Трансформаторы тока по сторонам ВН и НН собраны в <<звезду>>, полярными выводами к трансформатору.

\item Номинальный ток трансформатора со стороны ВН:
\begin{equation*}
    I_{\textsl{ном,в}} = \frac{S_{\textsl{ном}}}{\sqrt{3} \cdot U_{\textsl{ном,в}}}  = \frac{4 \cdot 10^6}{\sqrt{3} \cdot 37 \cdot 10^3} = 62,4\; (\textsl{А}),
\end{equation*}
\item Номинальный ток трансформатора со стороны НН:
\begin{equation*}
    I_{\textsl{ном,н}} = \frac{S_{\textsl{ном}}}{\sqrt{3} \cdot U_{\textsl{ном,н}}}  = \frac{4 \cdot 10^6}{\sqrt{3} \cdot 11 \cdot 10^3} = 210\; (\textsl{А}),
\end{equation*}

\end{enumerate}